% PAGES: 129-130

% Continuation of sec04_c1_1.tex (page 128 / folio 112): the \item below
% ("If lambda = -2, then (4.2) and (4.4)...") and the enclosing
% \begin{examplex} are the ones opened in that file. Do not reopen either.
We conclude that the eigenvalue $\lambda = -2$ has two linearly independent
eigenvectors, $\begin{pmatrix}1\\0\\2\end{pmatrix}$, and
$\begin{pmatrix}0\\1\\0\end{pmatrix}$.
\item If $\lambda = 1$, then (4.2) and (4.4) are both equivalent to $v_3 = 3v_1$, and
(4.3) is equivalent to $v_2 = -0.5v_1$. Thus, any eigenvector corresponding to
$\lambda = 1$ is of the form
\[
\begin{pmatrix}v_1\\ -0.5v_1\\ 3v_1\end{pmatrix} \;=\; 0.5v_1\begin{pmatrix}2\\-1\\6\end{pmatrix},
\]
and we conclude that the eigenvalue $\lambda = 1$ has one eigenvector, e.g.,
$\begin{pmatrix}2\\-1\\6\end{pmatrix}$.
\end{itemize}

(ii) Similarly, it can be shown that the matrix
\[
B \;=\; \begin{pmatrix} 2 & 2 & -1\\ -7 & -4 & 2.5\\ 2 & 4 & -1\end{pmatrix}
\]
has eigenvalues $1$ and $-2$, and has only one linearly independent eigenvector
corresponding to each eigenvalue, e.g., the eigenvector
$\begin{pmatrix}1\\-1\\2\end{pmatrix}$ corresponding to $-2$ and the eigenvector
$\begin{pmatrix}0\\1\\2\end{pmatrix}$ corresponding to $1$.
\end{examplex}

\begin{definition}
Let $A$ be a square matrix of size $n$. The characteristic polynomial $P_A(t)$ of the
matrix $A$ is the polynomial of degree $n$ given by
\begin{equation}
P_A(t) \;=\; \det(tI-A),
\end{equation}
where $I$ is the identity matrix of size $n$.
\end{definition}

\begin{theorem}
If $A$ is a square matrix, then $\lambda$ is an eigenvalue of $A$ if and only if
$\lambda$ is a root of the characteristic polynomial $P_A(t)$ of $A$. i.e., if and only
if
\[
\det(\lambda I - A) \;=\; 0.
\]
\end{theorem}

\begin{proof}
If $\lambda$ is an eigenvalue of $A$ with $v \neq 0$ a corresponding eigenvector, then
$Av = \lambda v$, which can be written as $(\lambda I - A)v = 0$. From (1.78) and
(1.81), it follows that
\begin{align*}
Av = \lambda v &\iff (\lambda I - A)v = 0 \iff v \in \Null(\lambda I - A)\\
&\iff \Null(\lambda I - A) \neq \{0\} \iff \lambda I - A \ \text{singular matrix}\\
&\iff \det(\lambda I - A) = 0.
\end{align*}
We conclude that $\lambda$ is an eigenvalue of $A$ if and only if
$P_A(\lambda) = \det(\lambda I - A) = 0$.
\end{proof}

The result below is a direct consequence of Theorem 4.1:

\begin{lemma}
Let $A$ be a square matrix of size $n$ and let $\lambda_i$, $i = 1:n$, be the
eigenvalues of $A$. The characteristic polynomial $P_A(t) = \det(tI-A)$ of $A$ can be
written as
\begin{equation}
P_A(t) \;=\; \prod_{i=1}^n (t-\lambda_i).
\end{equation}
\end{lemma}

% The "Examples:" (plural) discussion below revisits the two examples from the
% start of section 4.1 using characteristic polynomials. The source labels it
% "Examples:" (not "Example:"), so it is written out manually rather than via
% the \examplex environment (whose hard-coded label is singular) -- see the
% same convention used in ch02/sec02_c4_2.tex. This block is NOT closed in
% this file: it continues into sec04_c1_3.tex (page 131 / folio 115) and is
% closed there with a manual "\hfill$\square$\par\medskip" once part (ii)
% (matrix $B$) concludes. sec04_c1_3.tex must not print a second
% "Examples:" label for this block.
\par\medskip\noindent\textit{Examples:}\ We revisit the examples from the beginning of
this section, and identify the eigenvalues of the matrices by using characteristic
polynomials.

(i) Let
\[
A \;=\; \begin{pmatrix} -8 & 0 & 3\\ 3 & -2 & -1.5\\ -18 & 0 & 7\end{pmatrix}.
\]
Using formula (10.8) for computing the determinant of a $3 \times 3$ matrix, we find
that the characteristic polynomial of $A$ is
\begin{align*}
P_A(t)
&\;=\; \det(tI-A) \;=\;
\begin{vmatrix} t+8 & 0 & -3\\ -3 & t+2 & 1.5\\ 18 & 0 & t-7\end{vmatrix}\\
&\;=\; (t+8)(t+2)(t-7) - (-3)(t+2)(18) \;=\; (t+2)(t^2+t-2)\\
&\;=\; (t+2)^2(t-1).
\end{align*}

The roots of $P_A(t)$ are $\lambda_1 = \lambda_2 = -2$ and $\lambda_3 = 1$. We identify
the eigenvectors corresponding to each eigenvalue separately:

\begin{itemize}
\item If $\lambda = -2$, any corresponding eigenvector $v \neq 0$ is a solution to
$Av = -2v$, which can be written as
\[
\begin{cases}
-8v_1+3v_3 = -2v_1\\
3v_1-2v_2-1.5v_3 = -2v_2\\
-18v_1+7v_3 = -2v_3
\end{cases}
\quad\Longleftrightarrow\quad
\begin{cases}
v_3 = 2v_1\\
v_3 = 2v_1\\
v_3 = 2v_1
\end{cases}
\]
Thus, any eigenvector corresponding to the eigenvalue $\lambda = -2$ is of the form
\[
v \;=\; \begin{pmatrix}v_1\\ v_2\\ 2v_1\end{pmatrix} \;=\;
v_1\begin{pmatrix}1\\0\\2\end{pmatrix} + v_2\begin{pmatrix}0\\1\\0\end{pmatrix},
\]
and we conclude that the eigenvalue $\lambda = -2$ has two linearly independent
eigenvectors, e.g., $\begin{pmatrix}1\\0\\2\end{pmatrix}$ and
$\begin{pmatrix}0\\1\\0\end{pmatrix}$.
\item If $\lambda = 1$, any corresponding eigenvector $v \neq 0$ is a solution to
$Av = v$, which can be written as
\[
\begin{cases}
-8v_1+3v_3 = v_1\\
3v_1-2v_2-1.5v_3 = v_2\\
-18v_1+7v_3 = v_3
\end{cases}
\quad\Longleftrightarrow\quad
\begin{cases}
v_3 = 3v_1\\
3v_2 = 3v_1 - 1.5v_3\\
v_3 = 3v_1
\end{cases}
\]
\[
\Longleftrightarrow\quad
\begin{cases}
v_3 = 3v_1\\
v_2 = -0.5v_1\\
v_3 = 3v_1
\end{cases}
\]
% All display math above is complete (no \[ or \begin{cases} left open).
% CONTINUES in sec04_c1_3.tex (page 131 / folio 115): this \item, the
% enclosing \begin{itemize} (case analysis for matrix A revisited via the
% characteristic polynomial), and the manual "Examples:" block opened above
% are all left OPEN at the end of this file. sec04_c1_3.tex picks up
% directly with the conclusion of the lambda = 1 case for matrix A ("Thus,
% any eigenvector corresponding to the eigenvalue lambda = 1 is of the
% form..."), then closes this \itemize, gives part (ii) (matrix B), and
% closes the manual "Examples:" block with "\hfill$\square$\par\medskip"
% once matrix B's case analysis concludes (page 132 / folio 116).
